\documentclass[11pt]{article}
\usepackage[T1]{fontenc}
\usepackage{amsmath,amssymb,amsthm}
\usepackage[margin=1in]{geometry}
\usepackage{graphicx,microtype,xcolor}
\usepackage[hidelinks]{hyperref}
\newtheorem{theorem}{Theorem}
\newcommand{\R}{\mathbb R}
\title{Weighted dyadic summability does not control spectral-multiplier variation}
\author{Automated Math Discovery Run \texttt{fa\_banach\_001}}
\date{21 August 2026}
\begin{document}
\maketitle
\begin{center}
\fcolorbox{black}{red!5}{\parbox{0.88\linewidth}{
\textbf{Result.} The weighted replacement proposed after Theorem~3.3 of
Kriegler--Weis fails.  For \(c=4\) there is a bounded smooth multiplier
\(m\), with \(m(0)=0\), such that
\[
 \sum_{n\in\mathbb Z}\frac{\|m(2^n\cdot)\varphi_0\|_{W^4_2(\mathbb R)}}{1+|n|}<\infty,
\]
but, already for the identity operator on a one-point space,
\(t\mapsto m(tA)1=m(t)\) has infinite \(q\)-variation for every finite \(q\).
}}
\end{center}

\section{Source question}
The concluding remarks of \cite{KW} ask whether, in the \(q\)-variation
part of Theorem~3.3 and the time-dependent maximal estimate of Theorem~5.1,
the unweighted dyadic Sobolev sum can be relaxed by dividing its \(n\)-th
term by \(1+|n|\) (possibly with an additional logarithm).
\begin{figure}[ht]
\centering
\includegraphics[width=0.96\linewidth,height=0.48\textheight,keepaspectratio]{figures/open_question_crop.png}
\caption{The exact concluding questions in arXiv:2404.01893, PDF page 31.}
\end{figure}

\section{Idea of the proof}
The denominator \(1+|n|\) makes Sobolev cost cheap at remote dyadic
indices.  We place the \(j\)-th oscillatory block near the log-dyadic index
\(N_j=2^{6j}\), but give it only frequency \(K_j=2^j\).  Differentiating four
times costs \(O(K_j^4)\), so the proposed weighted sum is bounded by
\(C\sum_j K_j^4/N_j=C\sum_j2^{-2j}\).

Variation measures a different feature.  On the core of the \(j\)-th block,
the cutoff is identically one and the multiplier makes order \(K_j\)
alternating jumps of fixed size.  Thus its \(q\)-variation is at least a
constant times \(K_j^{1/q}\), which diverges.  Taking \(A=I\) on a one-point
space turns the operator conclusion into precisely this scalar variation
statement, while satisfying all functional-calculus hypotheses.

\section{Construction and proof}
Fix a nonzero \(\eta\in C_c^\infty((-1/4,1/4))\) such that
\(0\le\eta\le1\) and \(\eta=1\) on \([-1/8,1/8]\).  Put
\[
 N_j=2^{6j},\qquad K_j=2^j\quad(j\ge1),
\]
and, for \(\lambda>0\), define
\[
 m(\lambda)=\sum_{j\ge1}
 \eta(\log_2\lambda-N_j)
 \sin\bigl(2\pi K_j(\log_2\lambda-N_j)\bigr),
 \qquad m(0)=0.
\]
The log-supports are disjoint, so the sum is locally finite,
\(m\in C^\infty([0,\infty))\), and \(|m|\le1\).

\begin{theorem}
For every fixed dyadic cutoff
\(\varphi_0\in C_c^\infty((1/2,2))\),
\[
 \sum_{n\in\mathbb Z}
 \frac{\|m(2^n\cdot)\varphi_0\|_{W^4_2(\R)}}{1+|n|}<\infty,
 \qquad \|m\|_{V^q(0,\infty)}=\infty\quad(1\le q<\infty).
\]
\end{theorem}

\begin{proof}
On the support of \(\varphi_0\), the variable \(u=\log_2\lambda\) ranges
over a fixed compact interval.  Thus
\(m(2^n\lambda)\varphi_0(\lambda)\) can meet the \(j\)-th log-support only
when \(|n-N_j|\le C_0\), with \(C_0\) independent of \(j\); for each \(j\)
there are only \(O(1)\) such integers \(n\).  The chain rule on this fixed
compact \(\lambda\)-interval gives
\[
 \|m(2^n\cdot)\varphi_0\|_{W^4_2}
 \le C(1+K_j)^4\le C'K_j^4
\]
for those \(n\).  Since \(1+|n|\ge cN_j\) there,
\[
 \sum_n\frac{\|m(2^n\cdot)\varphi_0\|_{W^4_2}}{1+|n|}
 \le C\sum_{j\ge1}\frac{K_j^4}{N_j}
 =C\sum_{j\ge1}2^{-2j}<\infty.
\]

On the core interval \(|\log_2t-N_j|\le1/8\), the cutoff equals one.
Sampling successively at a maximum and the following minimum of the sine
wave produces at least \(K_j/4-2\) increments of magnitude \(2\).  Hence
\[
 \|m\|_{V^q}\ge2(K_j/4-2)^{1/q}.
\]
As \(K_j\to\infty\), this proves infinite \(q\)-variation.
\end{proof}

\section{Operator hypotheses and scope}
Let both underlying measure spaces be one point, let \(p=2\), \(Y=\mathbb C\),
and let \(A=I\).  Then \(A\) is \(0\)-sectorial, has a bounded
H\"ormander calculus of every admissible order, and \(e^{-tA}=e^{-t}I\) is
positive and contractive.  Taking \(\alpha=1\), the choice \(c=4\) satisfies
the smoothness inequality in Theorem~3.3.  For \(f=1\), however,
\(m(tA)f=m(t)\), so the asserted variational conclusion fails.

This disproves the proposed relaxation in the \(q\)-variation context and
therefore gives a negative answer to the blanket question.  It does not
settle whether the analogous weighted assumption might suffice for the
time-dependent maximal estimate, nor the separate square-summability
proposal.  The optional extra factor \(\log(|n|+2)\) in the numerator is
also harmless here, since \(\sum_j j\,2^{-2j}<\infty\).

\begin{thebibliography}{9}
\bibitem{KW} C.~Kriegler and L.~Weis,
\emph{\(q\)-variational H\"ormander functional calculus and Schr\"odinger
and wave maximal estimates}, arXiv:2404.01893.
\end{thebibliography}
\end{document}
